\documentclass[11pt]{article}

\usepackage[margin=1in]{geometry}
\usepackage[T1]{fontenc}
\usepackage[utf8]{inputenc}
\usepackage{amsmath,amssymb,array}
\usepackage[hidelinks]{hyperref}

\newcommand{\IUT}{IUT}
\newcommand{\Lean}{Lean}
\newcommand{\R}{\mathbb{R}}
\newcommand{\F}{\mathbb{F}}
\newcommand{\ThetaPilot}{\Theta\text{-pilot}}
\newcommand{\qPilot}{q\text{-pilot}}

\title{A First Lean Formalization Audit of the IUT III Corollary 3.12 Corridor}
\author{IUT Lean formalization project}
\date{May 29, 2026}

\begin{document}
\maketitle
\sloppy

\begin{abstract}
This note records the present state of a Lean formalization of a narrow part of
Mochizuki's inter-universal Teichmuller theory: the passage from IUT III,
Theorem 3.11 to Corollary 3.12, with attention to the Scholze--Stix objection.
The current Lean code proves conditional route theorems, obstruction lemmas, and
the final ordered-real algebra in Step (xi-f).  It does not prove Corollary 3.12
and does not refute it.  The main formal result so far is that a supplied
hull/log-volume \(q\)-to-\(\Theta\) comparison gives \(C_\Theta\ge-1\), and that a
nonarchimedean \((\mathrm{Ind3})\) entry together with explicit Hodge/SHE
square-weight transport data, or the primitive factored square/full-label SHE
obligations, feeds such a comparison, while several simplified collapse
mechanisms are rejected by Lean.
\end{abstract}

\section{Scope}

The review used the local Markdown conversions of IUT I--IV, Mochizuki's 2026
formalization note, and the Scholze--Stix note.  The relevant source passages are:
\begin{itemize}
\item IUT I: initial \(\Theta\)-data, prime strips, Hodge theaters, and
  non-scheme-theoretic \(\Theta\)-links.
\item IUT II: Hodge--Arakelov evaluation, theta values of shape \(q^{j^2}\),
  multiradiality, and Gaussian monoids.
\item IUT III: Theorem 3.11; Corollary 3.12, especially Steps (x) and (xi);
  Remarks 3.12.2 and 3.12.3.
\item IUT IV: later explicit log-volume estimates and Diophantine consequences.
\item Scholze--Stix: the ordered real-line identification problem in their
  Section 2.2.
\item Mochizuki 2026: the ``3.11.5 \(\Rightarrow\) 3.12'' fourth-triangle
  decomposition.
\end{itemize}
The namespace \texttt{Iut.Stage1} is local terminology for the present
corridor layer.  It is inspired by Mochizuki's staged formalization plan, but
does not assert completion of his Stage 1: many Frobenioid, Hodge-theater,
log-Kummer, and theta-link constructions remain external obligations.

\section{Mathematical Target}

IUT III Step (x) treats procession-normalized log-volumes as invariants for
\((\mathrm{Ind1})\) and \((\mathrm{Ind2})\), and turns \((\mathrm{Ind3})\) into
an upper inequality.  Step (xi) then uses the multiradial construction, IPL, SHE,
holomorphic hulls, determinants, and log-volume to place the \(q\)-pilot
log-volume in an upper ray controlled by the \(\Theta\)-pilot log-volume:
\[
  -|\log(q)| \in \R_{\le -|\log(\Theta)|}
  \quad\Rightarrow\quad
  -|\log(q)| \le -|\log(\Theta)|.
\]
If \(C_\Theta\) is any real number satisfying
\[
  -|\log(\Theta)|\le C_\Theta|\log(q)|,\qquad |\log(q)|>0,
\]
then Step (xi-f) concludes \(C_\Theta\ge-1\).  This last inference is now a
Lean theorem.  Lean also proves the strict variant: if the q-pilot lies strictly
inside the \(\Theta\)-upper ray, then \(C_\Theta>-1\); conversely, the boundary
case \(C_\Theta=-1\) forces equality of the q-pilot and \(\Theta\)-boundary
log-volumes and places the \(\Theta\)-boundary on the negative branch.
Mochizuki's 2026 note isolates the final comparison as:
\[
\begin{array}{ccccc}
\text{Theorem 3.11}
& \xrightarrow{\mathrm{APT}+\mathrm{IPL}}
& \text{``3.11.5''}
& \xrightarrow{\mathrm{SHE}+\mathrm{IPL}}
& \text{Corollary 3.12}.
\end{array}
\]
The present formalization targets only this corridor, not the full construction
of Theorem 3.11.

\section{Dispute Interface}

Scholze--Stix require explicit bookkeeping for all ordered one-dimensional real
vector spaces.  Their pressure point is that the \(\Theta\)-side concrete
objects carry \(j^2\)-scaled degrees, while a consistent identification of the
real lines appears to leave no room for inserting all these scalars.

The formalization therefore separates:
\[
\begin{array}{c}
\text{representative } j^2 \text{ data}\\
\text{balanced sign-compatible data}\\
\text{aggregate averaged data}\\
\text{final hull/log-volume data}
\end{array}
\qquad
\text{and}
\qquad
\begin{array}{c}
\text{raw real equality}\\
\text{transported equality}\\
\text{matched transport scales}\\
\text{cancellable ordered-line comparison.}
\end{array}
\]
This separation is the main correctness guard in the current code.

\section{Lean Coverage}

The relevant implementation is in
\texttt{Iut/Stage1/IUTStage1SourceCore.lean},
\texttt{Iut/Stage1/IUTStage1Remark312Absorption.lean},
\texttt{Iut/Stage1/IUTStage1Source.lean}, and
\texttt{Iut/Stage1/IUTStage1Experiments.lean}.
The code currently covers:
\begin{itemize}
\item labeled real-line copies and positive-scale transports;
\item a finite \(\F_\ell\)-label model with representative square weights
  \(j \mapsto j^2\);
\item a representative pilot-degree profile
  \(\deg(\Theta_j)=j^2\deg(q)\);
\item a finite \(q^{j^2}\) endpoint connecting this representative
  Gaussian-degree rule to the square-weight profile: for canonical
  representatives, the Gaussian degree is the square weight times the
  environment/q-pilot degree; Lean packages this together with the zero branch,
  the canonical \(j=1\) branch, and sign invariance;
\item a splitting-monoid/tensor-packet endpoint for the same finite branch:
  the core generator is zero, each procession generator is the procession
  square times the environment degree, acting on the tensor packet adds the
  generator log-volumes, and nonnegative environment degree gives monotonicity
  of the normalized acted packet; Lean now also identifies the averaged
  splitting-generator contribution with the full-label Gaussian average, so the
  normalized acted packet is the original packet plus either of these equivalent
  finite \(q^{j^2}\) averages; normalized-log-volume triviality of the action is
  equivalent in this finite model to zero environment degree, with strict
  increase for positive environment degree and strict decrease for negative
  environment degree;
\item proof that one label-independent scale cannot absorb all representative
  \(j^2\) scales in this model;
\item proof that, for nonzero \(\deg(q)\), one label-independent scale cannot
  account for all representative \(\Theta_j\)-degrees; Lean now proves the exact
  converse classifier: such a scale exists if and only if \(\deg(q)=0\); under
  the Corollary 3.12 sign hypothesis \(0<-\deg(q)\), no such scalar exists;
\item proof that the raw representative \(j^2\) theta-degree rule does not
  descend to the sign quotient \(|F_\ell|=F_\ell/\{\pm1\}\) when \(\deg(q)\ne0\);
\item an absolute-label exponent on \(|F_\ell|=\{0\}\cup F_\ell^\times/\{\pm1\}\)
  which agrees with \(j^2\) for the canonical representatives
  \(0\le j\le(\ell-1)/2\), and the corresponding theta-degree profile;
  Lean proves that the coordinate map \(F_\ell\to |F_\ell|\) is surjective,
  that the zero absolute label is reached exactly by \(j=0\), hence has
  singleton coordinate fiber, and that every nonzero sign-label class has a
  nonzero coordinate representative; Lean also
  proves the exact fiber relation:
  two coordinates have the same full absolute label if and only if they are
  equal or negatives of one another; equivalently, every nonzero full-label
  fiber is exactly a sign pair \(\{j,-j\}\), hence has cardinality \(2\).  These
  zero, canonical-one, equality-up-to-sign, and surjectivity facts are now
  packaged as a full-label endpoint.
  Lean also
  proves that the half-range procession map to \(|F_\ell|\) is bijective, using
  the minimal absolute representative \(j_{\min}\) for surjectivity and the
  exponent \(j^2\) for injectivity; this gives
  \(\#|F_\ell|=(\ell+1)/2\) in the finite model and permits Lean to reindex
  absolute-label averages by half-range processions;
  Lean also proves that the half-range procession top is exactly
  \((\ell-1)/2\) and that the resulting finite index set has cardinality
  \((\ell+1)/2\), matching the \(j\in\{0,\ldots,(\ell-1)/2\}\) convention in
  IUT II; on this half-range procession, Lean proves the Gaussian degree formula
  \(G(j)=j^2\deg(q)\), with \(G(0)=0\), and the normalized average identity
  \(6\,\operatorname{Avg}(G)=j_{\max}(2j_{\max}+1)\deg(q)\); by reindexing,
  the same identity holds for the average over \(|F_\ell|\).  Since
  \(j_{\max}\ge2\), Lean proves that this full average is not the singleton
  \(j=1\) identity value when \(\deg(q)\ne0\); in the nonpositive environment
  branch, the full \(|F_\ell|\)-average is bounded above by \(\deg(q)\), and
  hence by any bound above \(\deg(q)\); if \(\deg(q)<0\), this bound is strict:
  the full average is \(<\deg(q)\), while for \(\deg(q)>0\) it is \(>\deg(q)\);
  the same facts are now packaged as a
  typed \(|F_\ell|\)-indexed procession-normalized average object, where the
  average is provably distinct from the canonical \(j=1\) component unless
  \(\deg(q)=0\), strictly below it if \(\deg(q)<0\), strictly above it if
  \(\deg(q)>0\), and equal to that component if and only if \(\deg(q)=0\);
\item the degree-level Gaussian evaluation
  \(q\mapsto(q^{j^2})_{j\in |F_\ell|}\), including identity at \(j=1\);
  Lean proves that the absolute-label exponents at \(j=1\) and \(j=2\) are
  \(1\) and \(4\), hence the corresponding theta/Gaussian degrees are distinct
  whenever \(\deg(q)\ne0\), and equal exactly in the degenerate case
  \(\deg(q)=0\);
  Lean now also converts such an evaluation into a cusp/full-label log-volume
  compatibility object and proves the full-label value-preservation branch when
  the source and target environment degrees agree.  Conversely, Lean proves
  that full-label value preservation for two Gaussian evaluations is equivalent
  to equality of their environment degrees, and this equality is already
  detected at the canonical \(j=1\) component.  Lean also proves the IUT II
  restriction check at the canonical \(j=1\) sign label:
  the Gaussian component and the corresponding cusp-class log-volume are both
  the original environment degree; Lean also records the accompanying IUT II
  \(F_\ell\)-label model endpoint: the carrier has cardinality \(\ell\), the
  maps to and from \(\mathbb{Z}/\ell\mathbb{Z}\) are inverse, and the additive
  torsor action is transported by coordinate addition.  Lean also packages the
  unit/sign bridge: the distinguished sign unit acts by negation, the
  \(\{\pm1\}\)-unit subgroup orbits are exactly sign orbits, and the canonical
  cusp class is the sign class of coordinate \(1\).  It also records the
  zero/nonzero local-object assignment and the coordinate-to-full-label
  log-volume compatibility endpoint, including invariance under \(j\mapsto -j\);
  Lean now packages the Remark 3.12.2(i)(b) finite-symmetry split in one
  endpoint: additive \(F_\ell\)-translation, multiplicative unit action, the
  sign-unit orbit criterion, equality up to sign in the full-label quotient,
  and compatibility of unit action with coordinate-to-full-label descent,
  and the warning in the finite model: the singleton restriction to \(j=1\) is not
  stable under the additive \(F_\ell\)-torsor action, since translation by \(1\)
  sends \(1\) to \(2\ne1\); more generally, Lean proves that every nonempty
  translation-stable finite subset of \(F_\ell\) is all of \(F_\ell\), packaged
  as a translation-closure endpoint, so no nonempty proper subset restriction is
  compatible with this symmetry; Lean
  also proves that no nonzero additive translation descends to a well-defined
  map on \(|F_\ell|\), since the equal absolute labels of \(t\) and \(-t\) would
  have to map simultaneously to the nonzero label of \(2t\) and to \(0\);
  equivalently, an additive translation descends to \(|F_\ell|\) if and only if
  it is the zero translation; the same iff is exposed at the route interface as
  the exact condition for a translation equivalence to preserve full labels;
\item a canonical representative square-weight profile on \(F_\ell\):
  Lean proves the total \(j^2\)-mass is positive from the \(j=1\) contribution
  and therefore constructs the profile without an external positivity premise;
  it also proves
  \(6\sum_{j\in\mathbb{Z}/\ell\mathbb{Z}}j_{\mathrm{rep}}^2
   =(\ell-1)\ell(2(\ell-1)+1)\), hence the corresponding average formula after
  division by \(\ell\);
\item the square-weighted constant-average check:
  a constant procession-normalized log-volume remains equal to the same
  constant after averaging with the canonical \(j^2\)-weights;
\item the same constant-average check inside the
  \(F_\ell=\mathbb{Z}/\ell\mathbb{Z}\) cusp-label corridor:
  if the full-label log-volume is constant, the square-weighted average used in
  the \(\Theta\)-route is that constant, including for the canonical weights;
\item the constant-target structured-SHE subcase:
  the formal route now proves equality of the square-weighted average with the
  \(\Theta\)-source average, not merely an upper bound;
\item a log-volume shadow of the IUT III splitting-monoid action: the generator
  attached to a procession label contributes \(j^2\deg(q)\), contributes \(0\) at
  the core label, and adds its finite generator sum to the tensor-packet
  log-volume; Lean also proves the procession-normalized form: the acted
  tensor-packet average is the original tensor-packet average plus the averaged
  Gaussian generator contribution, equivalently the average of the
  \(j^2\deg(q)\) contributions over the procession container; for
  \(S_{j+1}=\{0,\ldots,j\}\), Lean expands the total generator contribution to
  \(j(j+1)(2j+1)\deg(q)/6\), and the resulting normalized log-volume change to
  \(j(2j+1)\deg(q)/6\); at the final procession length this averaged generator
  contribution is equal to the full-label Gaussian average; since
  \(j=\lfloor\ell/2\rfloor\ge2\), a nonnegative
  environment degree makes the acted normalized tensor-packet log-volume an
  upper bound for the original one; equality of acted and original normalized
  log-volume is equivalent to zero environment degree, and the sign of the
  environment degree determines the strict direction of the change;
\item the finite procession containers
  \(S^\pm_{j+1}=\{0,\ldots,j\}\), their nested inclusions, core-label stability,
  stage cardinality \(j+1\), and total procession ambiguity \((\ell^\pm)!\)
  rather than \((\ell^\pm)^{\ell^\pm}\);
\item the map from the final procession container to \(|F_\ell|\), sending the
  core label to \(0\) and recovering the exponent \(j^2\);
\item the log-volume shadow of the tensor packet over \(S^\pm_{j+1}\), with
  each factor a finite-fiber direct sum over places above \(v_\mathbb{Q}\), total
  log-volume given by the finite sum over the container, normalized denominator
  \(j+1\), and invariance under container permutations;
\item source-side local log-volume interfaces:
  Lean packages finite local log-volume identification, positive capsule count,
  procession-normalized averaging, and local container estimates
  \(q_{\mathrm{signed}}\le\operatorname{logvol}_{\mathrm{local}}\), while
  leaving the analytic construction of the local objects external;
\item the \(F_\ell\)-label average specialization:
  for labels represented as \(\mathbb{Z}/\ell\mathbb{Z}\), Lean rewrites the
  procession-normalized average denominator as \(\ell\);
\item the zero/nonzero Gaussian average split:
  Lean proves
  \[
  \#(\mathbb{Z}/\ell\mathbb{Z})^\times=\ell-1,\qquad
  \operatorname{Avg}_{j\in F_\ell}G(j)
  =\frac{\ell-1}{\ell}\operatorname{Avg}_{j\ne0}G(j),
  \]
  for the degree-level Gaussian evaluation \(G\), using \(G(0)=0\).  Thus the
  paper's uniform \(1/\ell\) weighted-volume convention is not silently replaced
  by the auxiliary nonzero-only average; in the negative signed branch Lean also
  proves the strict inequality
  \(\operatorname{Avg}_{j\ne0}G(j)<\operatorname{Avg}_{j\in F_\ell}G(j)\),
  while in the positive branch the inequality is reversed; equality between the
  two averages holds if and only if \(\deg(q)=0\).  Lean also proves that the
  all-\(F_\ell\) coordinate average itself is zero if and only if
  \(\deg(q)=0\);
  Lean separately formalizes the nonzero half-range
  \(j=1,\ldots,(\ell-1)/2\) average, proves that the signed nonzero carrier
  \(F_\ell^\times\) is the twofold cover of this half-range by \(j\) and \(-j\),
  and therefore proves that the signed nonzero average equals the nonzero
  absolute-label average.  It also proves
  \(6\,\operatorname{Avg}_{j\ne0,\,|F_\ell|}G(j)
   =(j_{\max}+1)(2j_{\max}+1)\deg(q)\); Lean also proves the exact rescaling
  \[
    \operatorname{Avg}_{|F_\ell|}G
    =\frac{j_{\max}}{j_{\max}+1}\operatorname{Avg}_{j\ne0,\,|F_\ell|}G.
  \]
  Since the full \(|F_\ell|\)-average
  has coefficient \(j_{\max}(2j_{\max}+1)/6\), Lean proves that the nonzero
  absolute-label average is strictly below the full average when \(\deg(q)<0\)
  and strictly above it when \(\deg(q)>0\), with equality if and only if
  \(\deg(q)=0\).  Combining the signed and absolute computations, Lean proves
  \[
  \operatorname{Avg}_{F_\ell}G
  =\frac{j_{\max}(j_{\max}+1)}{3}\deg(q),
  \]
  equivalently
  \(\operatorname{Avg}_{F_\ell}G=((\ell+1)/\ell)
  \operatorname{Avg}_{|F_\ell|}G\),
  so \(\operatorname{Avg}_{F_\ell}G\) is strictly below
  \(\operatorname{Avg}_{|F_\ell|}G\) when \(\deg(q)<0\), strictly above it when
  \(\deg(q)>0\), and equal to it only when \(\deg(q)=0\);
\item the IUT II Remark 4.7.3 zero/nonzero weighted-volume relation:
  Lean records a relation in which every nonzero full label is subordinate to
  the zero label, proves that \(0\ll0\) is false, proves that this relation
  is not symmetric at the zero label, and proves the exact criterion
  \(a\ll0\iff a\ne0\) for full labels; at coordinate level this becomes
  \(j\ll0\iff j\ne0\), so the all-coordinate statement is false precisely
  because it includes \(j=0\).  Lean packages this with unit-action
  preservation and the zero-translation criterion for affine preservation.
  Lean also proves that the full-label set
  subordinate to zero has cardinality \((\ell-1)/2\), the nonzero absolute
  half-range, and decomposes any full-label sum, hence any full-label average,
  into its zero contribution plus the subordinate nonzero contribution with
  denominator \((\ell+1)/2\).  Applying this to the Gaussian degree function,
  whose zero component is \(0\), Lean rewrites the full absolute-label
  Gaussian average as the subordinate nonzero Gaussian sum divided by the full
  absolute-label denominator and proves the closed form
  \[
    6\sum_{a\ll0}G(a)
    =j_{\max}(j_{\max}+1)(2j_{\max}+1)\deg(q).
  \]
  Lean also proves that this subordinate sum is exactly the canonical
  nonzero absolute half-range sum \(\sum_{j=1}^{j_{\max}}G(j)\).
  In contrast to nonzero additive translations, Lean defines the multiplicative
  unit action on full labels, proves its compatibility with the coordinate
  formula \(j\mapsto a\cdot j\), and proves that it preserves the subordinate
  zero/nonzero relation.  Lean also proves the unit and multiplication laws for
  this action and proves that the sign subgroup \(\{\pm1\}\) acts trivially on
  full absolute labels.  Each unit action is packaged as an equivalence of
  full labels with inverse action by \(a^{-1}\), and Lean proves that the zero
  label is fixed exactly.  Consequently Lean proves that full-label sums and
  full-label averages are invariant under this multiplicative permutation; the
  same holds for sums over the subordinate nonzero full-label subset.  These
  invariance statements are also specialized to the Gaussian degree function
  \(G(a)\), so the full Gaussian average and the subordinate Gaussian sum are
  invariant under multiplicative reindexing.  At the coordinate level, Lean
  proves the corresponding sharp contrast: unit multiplication preserves the
  predicate \(j\ll0\), whereas additive translation preserves this predicate
  for all \(j\) if and only if the translation is zero.  More strongly, an
  affine coordinate operation \(j\mapsto a j+t\), with \(a\in F_\ell^\times\),
  descends to full absolute labels if and only if \(t=0\), and preserves the
  subordinate predicate for all \(j\) if and only if \(t=0\).  Lean now packages
  the unit-action laws, coordinate compatibility, zero-label preservation, and
  affine descent criterion in one endpoint.
  Lean separately proves that raw \(F_\ell\)-coordinate sums are invariant
  under additive translation, unit multiplication, and their affine
  composition; this is only finite-permutation invariance and does not supply
  descent to \(|F_\ell|\).  For the Gaussian coordinate average, Lean packages
  this as a diagnostic: every affine map with nonzero additive offset preserves
  the raw coordinate average but fails full-label descent.  If the Gaussian
  environment degree is nonzero, the same affine operation also changes the
  zero Gaussian component, so aggregate average invariance is not pointwise
  Gaussian compatibility.  Lean also proves the exact degeneracy criterion:
  a coordinate Gaussian component is zero if and only if its coordinate is zero
  or the environment degree is zero; for an affine zero-coordinate image this
  becomes \(t=0\) or \(\deg(q)=0\).
  At the purely labelwise level, Lean proves the exact pointwise preservation
  criterion: \(j\mapsto a j+t\) preserves the full absolute label of every
  coordinate if and only if \(t=0\) and \(a\in\{\pm1\}\).
  At the Gaussian-value level, Lean proves the exact nondegenerate pointwise
  criterion: if the environment degree is nonzero, then
  \(j\mapsto a j+t\) preserves all coordinate Gaussian degrees if and only if
  \(t=0\) and \(a\in\{\pm1\}\).  Equivalently, in this finite affine family,
  nondegenerate pointwise Gaussian preservation is exactly preservation of the
  full absolute-label map.  Lean also proves the degenerate complement:
  if \(\deg(q)=0\), then every Gaussian degree is \(0\), so every affine
  coordinate operation preserves Gaussian values.  Combining both cases, Lean
  proves the exact classification: affine Gaussian-value preservation is
  equivalent to either \(\deg(q)=0\) or preservation of the full absolute-label
  map.  Uniform preservation over all nonzero environment degrees is therefore
  equivalent exactly to full absolute-label preservation.
  For the stricter representative \(j^2\) square-weight branch, Lean proves the
  corresponding affine criterion \(t=0\) and \(a=1\).  Thus the sign action is
  legitimate for the balanced/full-label Gaussian branch, but not for the
  representative-square transport branch.  Requiring both the
  representative-square branch and the full-label branch is therefore
  equivalent, in the unit-affine family, to the identity case \(t=0,a=1\).
  Likewise, nondegenerate pointwise Gaussian preservation plus representative
  square preservation is equivalent to \(t=0,a=1\).
  Lean exposes this as a direct
  diagnostic: the affine action \(j\mapsto -j\) preserves Gaussian degrees
  pointwise but does not preserve the representative \(j^2\) coordinate-square
  profile.  Lean now packages the sign-negation case as a source-layer
  transport theorem:
  the balanced self-transport \(j\mapsto -j\) simultaneously preserves all
  Gaussian degrees and all balanced square weights, while it fails preservation
  of the representative constant-one square-weighted summands.  This is a
  formal separation between sign-compatible balanced transport and the stricter
  representative branch.  A further experiment theorem combines this
  preservation/failure package with the final-route audit: the same balanced
  negation transport is not the comparison level consumed by the final
  weighted-theta route.
  Lean also converts the
  zero/nonzero cusp-label compatibility object to a uniform \(F_\ell\) average
  and proves that a diagonal constant log-volume distribution averages to that
  same constant; the same compatibility layer now exposes, as a single endpoint,
  the equality between coordinate-normalized log-volumes and full-label
  log-volumes, together with the zero branch, nonzero cusp-class branch, and
  sign-pair invariance; Lean also proves that the uniform \(F_\ell\)-average is
  the average of these full-label values, and that zero/cusp-class upper bounds
  propagate to this average;
\item the IUT III Step (x) averaged-log-volume corridor:
  Lean records that \((\mathrm{Ind1})\) and \((\mathrm{Ind2})\) preserve the
  procession-normalized averaged log-volume, while \((\mathrm{Ind3})\) supplies
  an upper bound; the code now constructs this corridor from three
  \(F_\ell\)-cusp log-volume compatibility objects when the after-\((\mathrm{Ind2})\)
  zero and cusp-class branches are bounded by the proposed \((\mathrm{Ind3})\)
  bound; the before-, after-\((\mathrm{Ind1})\)-, and
  after-\((\mathrm{Ind2})\)-averages are all bounded by the same
  \((\mathrm{Ind3})\) upper bound; Lean now packages the
  Remark 3.12.2(iii) log-link-juggling endpoint by combining these two
  equalities with the one-sided \((\mathrm{Ind3})\) bound, and separately
  records the matching nonarchimedean/archimedean place families for the
  upper-semi state;
\item the Step (x) to Step (xi) upper-ray handoff:
  if the q-pilot reading is the pre-indeterminacy averaged log-volume and the
  \(\Theta\)-hull reading is the \((\mathrm{Ind3})\) upper bound, Lean constructs
  the hull/determinant upper-ray datum used in Step (xi); the code now also
  constructs this handoff directly from the \(F_\ell\)-cusp compatibility
  corridor plus the determinant/hull identifications and proves the resulting
  q-pilot membership in the hull upper ray; Lean now exposes the
  defining equivalence \(x\in R_{\le\Theta}\iff x\le\Theta\), so the
  Step (xi-f) passage from q-pilot membership in the upper ray to the ordered
  real inequality is explicit; the same handoff also constructs the finite
  extended \(\Theta\)-log-volume endpoint, excludes \(+\infty\), identifies its
  real value with the hull boundary, and proves the q-pilot comparison against
  this real value before adjoining any \(C_\Theta\)-hypothesis; Lean also
  identifies this finite endpoint with the signed two-computation endpoint:
  the signed \(q\)-reading is the finite endpoint's q-pilot log-volume, the
  signed \(\Theta\)-reading is its real \(\Theta\)-value, and the formal
  \texttt{Corollary312Inequality} is obtained from these two shared readings;
\item the Step (x) to Step (xi-g) two-computation handoff:
  the same data constructs the q-pilot two-computation record, with the input
  prime-strip reading equal to the pre-indeterminacy average and bounded by the
  \(\Theta\)-hull upper-ray value; Lean now also proves, at this route level,
  that the input prime-strip q-reading and the hull-side q-reading are equal,
  and hence that the pre-indeterminacy average is the hull-side q-reading; the
  Fig. 3.8 endpoint now packages both readings, their equality, both
  upper-ray inequalities, and hull-side q-membership in one theorem;
  before any \(C_\Theta\) is chosen, the resulting signed two-computation
  endpoint already yields the formal \texttt{Corollary312Inequality}, with
  q- and \(\Theta\)-pilot magnitudes identified as the negatives of the
  pre-indeterminacy average and the hull boundary;
\item the Step (x)-sourced \(C_\Theta\) endpoint:
  adding \(|\log(q)|>0\) for the pre-indeterminacy q-pilot average and the
  displayed upper bound
  \(-|\log(\Theta)|\le C_\Theta|\log(q)|\) yields \(C_\Theta\ge-1\);
  Lean also exposes the route-level chain
  \(q_{\mathrm{avg}}\le\Theta_{\mathrm{real}}\le C_\Theta|q_{\mathrm{avg}}|\),
  hence \(q_{\mathrm{avg}}\le C_\Theta|q_{\mathrm{avg}}|\), before invoking the
  final ordered-real algebra; the Step (xi-f) endpoint now packages q-pilot
  upper-ray membership, both ordered inequalities, \(C_\Theta\ge-1\), and the
  rejection of \(C_\Theta<-1\) directly at the route level; if the q-pilot
  inequality is strict, Lean proves \(C_\Theta>-1\); Lean now also
  combines this Step (xi-f) package with the Remark 3.12.2 distinct-label
  intertwining discipline, so the inequality endpoint is available while the
  q-native and \(\Theta\)-native labels remain unequal;
\item the Step (x)-sourced statement endpoint:
  after adjoining the statement-level
  \(\Theta\)-subject/\(q\)-not-subject boundary, the same data constructs the
  Corollary 3.12 statement endpoint, with finite \(\Theta\)-branch excluding
  \(+\infty\) and \(C_\Theta\ge-1\); Lean proves that this constructed endpoint still records
  the \(\Theta\)-pilot as subject to \((\mathrm{Ind1}),(\mathrm{Ind2}),
  (\mathrm{Ind3})\) and the \(q\)-pilot as not subject to them, and that these
  two status fields are unequal at the constructed endpoint; the same endpoint
  is now constructed directly from the \(F_\ell\)-cusp compatibility corridor,
  determinant/hull identifications, q-pilot positivity, and the displayed
  \(C_\Theta\)-bound, preserving the q-pilot average, upper-ray membership,
  real \(\Theta\)-value, and status separation; from the same \(F_\ell\)-sourced
  endpoint, Lean proves the nonnegative-\(\Theta\) branch:
  if the \((\mathrm{Ind3})\) upper bound is nonnegative, then the q-pilot
  average is strictly smaller than it and the endpoint has \(C_\Theta\ge0\);
  Lean also packages this branch as a formal `Corollary312Inequality` value
  for the comparison data induced by the endpoint sign reduction; without the
  nonnegative-\(\Theta\) assumption, the same \(F_\ell\)-sourced endpoint proves
  \(q_{\mathrm{avg}}\le C_\Theta|q_{\mathrm{avg}}|\), \(C_\Theta\ge-1\), and
  rejection of \(C_\Theta<-1\); the fixed-q two-computation version is also
  exposed, with the lower-bound algebra applied to
  \(-|\log(q)|\) as determined by the input prime-strip q-computation; the
  signed \(C_\Theta\) comparison payload is exposed from the same data, with
  signed \(q\) equal to the pre-indeterminacy average and signed \(\Theta\)
  equal to the hull boundary, and Lean derives the corresponding formal
  \texttt{Corollary312Inequality} from this signed payload while identifying
  the positive q- and \(\Theta\)-pilot magnitudes as the negatives of these
  two signed readings;
  Lean also proves
  that the endpoint \(q\)-pilot log-volume is exactly the pre-indeterminacy
  average and, by Step (x) invariance, also the after-\((\mathrm{Ind1})\)- and
  after-\((\mathrm{Ind2})\)-averages; thus \(|\log(q)|\) is the negative of
  these equal averages and is positive; equivalently, the three equal averages
  are negative under the q-pilot positivity hypothesis; Lean also proves
  the displayed inequality
  \(-|\log(q)|\le-|\log(\Theta)|\) at this endpoint, with \(q\) read as the
  pre-indeterminacy average and \(\Theta\) read as the finite real branch;
  equivalently, the endpoint \(q\)-pilot log-volume lies in the constructed
  upper ray \(\mathbb{R}_{\le-|\log(\Theta)|}\);
  Lean also formalizes the opening reduction in the printed proof: if this
  finite real \(\Theta\)-value is nonnegative, then the strict negativity of the
  q-pilot log-volume gives the comparison immediately; moreover the displayed
  bound
  \(-|\log(\Theta)|\le C_\Theta|\log(q)|\) then forces \(C_\Theta\ge0\);
  the finite extended \(\Theta\)-value and the real \(\Theta\)-branch are
  identified with the Step (x) \((\mathrm{Ind3})\) upper bound; the same endpoint
  bound holds after the
  \((\mathrm{Ind1})\) and \((\mathrm{Ind2})\) average-preserving
  transformations; the Step (x)-sourced endpoint also directly rejects
  \(C_\Theta<-1\);
\item base-valuation products of tensor packets, the finite
  \((F_{\mathrm{mod}}^\times)_j\)-action shadow, and a coric transfer shadow
  \(F\to D\to D^\vdash\) preserving product log-volume without identifying
  Hodge-theater histories; Lean now also packages the Remark 3.12.2(i)(c)
  ring-structure split: additive/log-shell direct sums, multiplicative/global
  Frobenioid calibration, the local-shift criterion, and the guarded transfer
  between histories;
\item the Remark 3.9.5 holomorphic-hull operation:
  Lean now represents hull formation by a closure operator on regions, proving
  the fixed-hull, extensivity, monotonicity, idempotence, and log-volume
  monotonicity properties P1--P3 used in Step (xi); from a q-region contained
  in the theta hull, Lean constructs the existing hull/determinant upper-ray
  endpoint \(q_{\mathrm{pilot}}\le\Theta_{\mathrm{hull}}\); Lean now also
  packages the Step (xi) union of possible theta images and derives the same
  upper-ray endpoint from containment of the q-pilot region in the holomorphic
  hull of that union; this possible-image hull endpoint now feeds directly into
  the \(C_\Theta\)-lower-bound algebra, giving \(C_\Theta\ge-1\), the
  boundary/strict split, and equality of the theta hull with the q-pilot
  log-volume in the boundary case; Lean also records the hull minimality
  principle: any closed region containing all possible theta images contains
  the theta hull, hence bounds its log-volume; the possible-image hull endpoint
  now also proves the strict branch \(C_\Theta>-1\) whenever the q-pilot
  log-volume is strictly below the theta hull, and packages the resulting
  boundary-equality-versus-strict alternative;
\item the Step (xi-d)--(xi-g) shadow: determinant normalization, upper-ray
  membership, and equality of the two \(q\)-pilot log-volume computations; the
  determinant normalization is invariant under changing the positive
  tensor-power representative when the determinant log-volume is fixed; equality
  of the \(q\)-pilot log-volume with the \(\Theta\)-hull or determinant boundary
  is equivalent to supplying the corresponding reverse inequality;
\item the Corollary 3.12 finiteness endpoint:
  the statement-level value \(-|\log(\Theta)|\in\mathbb{R}\cup\{+\infty\}\)
  is represented by an explicit finite/\(+\infty\) split, and the
  hull/determinant upper-ray datum forces the finite real branch, explicitly
  excluding the \(+\infty\) branch; Lean also exposes the resulting real
  \(\Theta\)-log-volume and proves the \(q\)-pilot upper-ray comparison against
  that real value;
\item the Step (xi-d) determinant normalization:
  the hull boundary is tied to a rank \(>1\) arithmetic-vector-bundle
  determinant, a positive tensor power of that determinant, and the normalized
  determinant log-volume; Lean now proves at the corridor endpoint that the
  normalized value is exactly the determinant log-volume itself, so the
  \(\Theta\)-hull log-volume is the determinant log-volume; it also records the
  equivalent tensor-power equations before and after normalization, and transports
  the q-pilot upper-ray comparison to both determinant coordinates; the same
  determinant/tensor-power description is now available at the finite
  \(\Theta\)-endpoint;
\item the Corollary 3.12 statement boundary:
  the \(\Theta\)-pilot log-volume is typed as subject to
  \((\mathrm{Ind1}),(\mathrm{Ind2}),(\mathrm{Ind3})\), while the \(q\)-pilot
  log-volume is typed as not subject to these indeterminacies;
\item the extended statement endpoint:
  from the finite \(\Theta\)-branch, \(0<-|\log(q)|\), the upper-ray comparison,
  and the displayed hypothesis
  \(-|\log(\Theta)|\le C_\Theta|\log(q)|\), Lean proves \(C_\Theta\ge-1\);
  if the finite \(\Theta\)-value is nonnegative, the q-pilot is strictly below
  it and Lean proves the sharper \(C_\Theta>-1\);
\item the Step (xi-f) final algebra:
  from \(-|\log(q)|\le-|\log(\Theta)|\), \(|\log(q)|>0\), and
  \(-|\log(\Theta)|\le C_\Theta|\log(q)|\), Lean proves \(C_\Theta\ge-1\);
  from the strict inequality
  \(-|\log(q)|<-|\log(\Theta)|\), Lean proves \(C_\Theta>-1\); if
  \(C_\Theta=-1\), Lean proves
  \(-|\log(q)|=-|\log(\Theta)|\), hence equality with the determinant
  log-volume, and \(-|\log(\Theta)|<0\); hence the endpoint splits into the
  boundary branch \(C_\Theta=-1\) or the strict branch \(C_\Theta>-1\); Lean also
  proves the determinant-coordinate strict branch: if the fixed q-reading is
  strictly below the determinant boundary, then \(C_\Theta>-1\); equivalently, at
  the determinant-coordinate endpoint, Lean separates the boundary branch where
  both fixed q-readings equal the determinant log-volume from the strict branch;
  combined with the tensor-power warning, the boundary branch also excludes both
  fixed q-readings from the \(N\)-tensor-power upper ray;
\item the signed-payload bridge:
  the existing Stage 1 comparison record \(q_{\mathrm{signed}}\le
  \Theta_{\mathrm{signed}}\), together with
  \(\Theta_{\mathrm{signed}}\le C_\Theta(-q_{\mathrm{signed}})\), feeds this
  Step (xi-f) algebra directly;
\item the Step (xi-g) endpoint bridge:
  the two-computation q-pilot record produces the signed comparison payload
  with \(q_{\mathrm{signed}}\) read from the input prime-strip log-volume and
  \(\Theta_{\mathrm{signed}}\) read from the hull/determinant upper-ray
  boundary; Lean now propagates the strict and boundary \(C_\Theta\)-branches to
  both q-computations: a strict hull-output inequality gives \(C_\Theta>-1\),
  while \(C_\Theta=-1\) forces the hull-output q-reading to equal the
  \(\Theta\)-boundary;
\item the Step (xi-e)--(xi-f) fixed-value condition:
  in the signed two-computation endpoint, Lean proves that the input
  prime-strip reading and the hull-output reading are both the same real value
  \(-|\log(q)|\), and that this value lies in the upper ray
  \(\mathbb{R}_{\le-|\log(\Theta)|}\); the source endpoint now also records these
  two fixed readings in determinant and normalized tensor-power coordinates, with
  equality to the determinant boundary guarded by the corresponding reverse
  determinant inequality;
\item the fixed-value \(C_\Theta\) endpoint:
  combining the preceding fixed upper-ray membership with
  \(-|\log(\Theta)|\le C_\Theta|\log(q)|\), Lean derives
  \(-|\log(q)|\le C_\Theta|\log(q)|\) and hence \(C_\Theta\ge-1\);
  consequently the same hypotheses reject any \(C_\Theta<-1\);
\item the statement-endpoint composition:
  the two-computation \(C_\Theta\) endpoint canonically supplies the finite
  \(\Theta\)-branch Corollary 3.12 statement endpoint once the
  \(\Theta\)-subject/\(q\)-not-subject indeterminacy boundary is provided;
\item the opening sign reduction in the proof of Corollary 3.12:
  from \(0<-q_{\mathrm{signed}}\), Lean proves that
  \(0\le\Theta_{\mathrm{signed}}\) immediately implies
  \(q_{\mathrm{signed}}\le\Theta_{\mathrm{signed}}\), matching the printed
  reduction to the case \(-|\log(\Theta)|<0\); the same strict-negativity
  argument is now available at the statement endpoint for the real
  \(\Theta\)-log-volume, and the statement endpoint now maps directly to the
  sign-reduction record and proves the nonnegative branch of the Corollary 3.12
  comparison from it;
\item the Step (xi-h) warning: for \(N\ge2\) and \(\Theta<0\), the tensor-power
  boundary satisfies \(N\Theta<\Theta\); Lean now also proves that the original
  \(\Theta\)-boundary lies in the original upper ray but not in the tensor-power
  upper ray, so the replacement changes the comparison region; at the
  \(F_\ell\)-sourced corridor endpoint, this warning is tied to the finite
  \(\Theta\)-branch: the \(N\)-tensor-power upper ray is a strict tightening of
  the finite endpoint's upper ray and excludes the original real \(\Theta\)
  value; Lean now also proves that every value in the tensor-power upper ray is
  strictly below the original \(\Theta\)-boundary; consequently, if the fixed
  input or output q-reading is forced into this tensor-power upper ray while the usual
  \(C_\Theta\)-bound is assumed, Lean derives the strict branch
  \(C_\Theta>-1\), hence the boundary branch \(C_\Theta=-1\) forbids such
  q-membership in the tensor-power upper ray; the same exclusion now applies to
  both q-readings of the two-computation endpoint, and Lean packages the
  resulting endpoint alternative as boundary-with-exclusion versus strict branch;
\item the Step (xii) shadow: local Frobenioid \(p_v^N\)-ambiguity changes
  log-volume by an integer shift, while global calibration fixes exponent \(0\);
  Lean now records the exact criterion: the local shifted and unshifted
  log-volumes agree, and likewise the globally calibrated and local shifted
  values agree, exactly when the local exponent-step product is zero, equivalently
  when the local exponent is zero or the local prime-step log-volume is zero; for
  a fixed local base and nonzero local step, equality of two shifted readings is
  equivalent to equality of their integer exponents; at the
  \(F_\ell\)-sourced finite \(\Theta\)-endpoint, the globally calibrated
  log-volume is identified with the finite real \(\Theta\)-value, while any
  nonzero local exponent with nonzero local step gives a distinct shifted
  local log-volume; Lean now also couples this global calibrated value to the
  same \(C_\Theta\)-endpoint that proves \(C_\Theta\ge-1\), and at that endpoint
  keeps the equality criterion in the decomposed exponent-or-step form; there Lean also
  records the ordered alternative: a positive exponent-step product places the
  local shifted log-volume above the calibrated value, and a negative product
  places it below; in the boundary branch \(C_\Theta=-1\), Lean now proves that
  the globally calibrated value agrees with the q/determinant value while any
  nonzero local Frobenioid shift remains distinct from that value; conversely,
  identifying the shifted local value with the boundary q/determinant value is
  equivalent to a trivial local shift, and the sign of the local exponent-step
  product determines whether the shifted local value lies above or below the
  boundary value; if the local prime-step log-volume is nonzero, this boundary
  identification is equivalent simply to exponent \(N=0\); if the local step is
  set to zero, every exponent collapses to the same boundary q/determinant value;
  the same ordered local-shift comparison is now stated against both fixed
  q-readings;
  Lean now fuses these facts into a trichotomy: at \(C_\Theta=-1\), the sign of
  the Frobenioid shift term decides whether the local shifted value is below,
  equal to, or above both fixed q-readings;
  Lean now also packages the determinant boundary, tensor-ray exclusion, and
  nonzero local-shift separation into one boundary-versus-strict branch theorem,
  including separation from both fixed q-readings of the two-computation endpoint;
  a parallel branch theorem packages the determinant/tensor boundary together
  with the ordered local-shift comparison around both fixed q-readings;
  this has now been sharpened to a determinant/tensor branch theorem carrying
  the full sign trichotomy for the local shifted value, opposed to the strict
  branch \(C_\Theta>-1\);
  with nonzero local step, equality with either fixed q-reading is equivalent to
  exponent \(N=0\), while in the zero-step case Lean proves the complementary
  collapse to both fixed q-readings;
\item the Remark 3.12.1(iii) arithmetic fact
  \(\mathbb{Q}_{>0}\cap\mathbb{Z}^{\times}=\{1\}\), formalized as the exact
  characterization \(0<q\iff q=1\) for rational integer units \(q=\pm1\);
\item the Remark 3.12.1(ii) generalized-link numerical bound
  \(C_\Theta\ge-\lambda\) for \(\lambda\in\mathbb{Q}_{>0}\), with strict
  strengthening over the \(\lambda=1\) bound when \(\lambda<1\); Lean now proves
  the resulting strict conclusion \(-1<C_\Theta\) in this case, and now also
  splits the generalized endpoint into the boundary branch
  \(C_\Theta=-\lambda\), where the \(q^\lambda\)-pilot and \(\Theta\)-pilot
  signed log-volumes are equal, or the strict branch \(C_\Theta>-\lambda\);
\item the \(q^\lambda\)-pilot endpoint algebra:
  from a signed comparison for \(-\lambda|\log(q)|\) and the same
  \(C_\Theta|\log(q)|\) upper bound, Lean derives \(C_\Theta\ge-\lambda\), and
  hence \(C_\Theta\ge-1\) when \(\lambda\le1\), with the strict
  \(-1<C_\Theta\) conclusion when \(\lambda<1\); the \(F_\ell\)-sourced
  corridor exposes this as a conditional generalized endpoint, without
  asserting that the ordinary \(\lambda=1\) construction supplies all
  generalized links; Lean now also specializes the endpoint at \(\lambda=1\)
  and proves that its signed \(q^\lambda\)-boundary is exactly the ordinary
  q-pilot log-volume; the strict and boundary branch classifiers specialize to
  the ordinary q-pilot upper-ray comparison at \(\lambda=1\);
\item the Remark 3.12.2 distinct-label implication
  \(q\text{-itw}\Rightarrow q\text{-itw}\wedge\Theta\text{-itw}/\mathrm{indets}\),
  with q-native and \(\Theta\)-native labels proved distinct; Lean now also
  packages this simultaneous validity together with the fixed weakened
  \(F^{\times\mu}\)-prime-strip condition, and exposes the second implication
  \(q\text{-itw}\wedge\Theta\text{-itw}/\mathrm{indets}
  \Rightarrow\Theta\text{-itw}/\mathrm{indets}\) as formal projection; at the
  \(F_\ell\)-sourced finite/signed endpoint, Lean combines the signed
  q-to-\(\Theta\) comparison with this distinct-label discipline, so the
  comparison is available while the q-native and \(\Theta\)-native labels remain
  provably unequal; the same discipline is now attached to the Step (xi-f)
  \(C_\Theta\)-endpoint itself;
\item the Remark 3.12.2 toy calculation: forgotten concrete assignments
  \(-h\) and \(-2h\) are incompatible for \(h>0\), while
  \(-h\in\R_{\le-2h+\epsilon}\) implies \(h\le\epsilon\); Lean now packages
  this as the full \((\mathrm{ftoy})\) endpoint:
  membership in the upper ray, the displayed inequality
  \(-h\le-2h+\epsilon\), and the bound \(h\le\epsilon\); the same toy endpoint
  is now instantiated from the \(F_\ell\)-sourced finite \(\Theta\)-upper ray
  whenever the corridor q-boundary and \(\Theta\)-boundary are written as
  \(-h\) and \(-2h+\epsilon\);
\item the Fig. 3.9 column table: Lean packages the upper-row similarities
  shared by the 0-column and 1-column, namely the essential
  \(\bullet_0\) and \(\circ\) roles, log-volume compatibility, and
  log-Kummer non-interference; the simultaneous Remark 3.12.2(iv) endpoint
  now records these facts for both columns, including q-pilot
  non-interference; Lean now also packages the q-pilot-specific combination:
  non-interference holds, multiradiality fails, and the log-shell treatment is
  tautological; it also packages the lower-row differences:
  \(\Theta\)-pilot multiradiality holds, q-pilot multiradiality does not, and
  the log-shell treatments differ;
\item the Remark 3.9.5(v) upper-semi set quotient:
  Lean models the quotient \(E^S\) that collapses a subset \(S\subseteq E\)
  to a single point and leaves the complement unchanged; every nonempty
  subset of \(S\) has image equal to the collapsed singleton, while outside
  points remain injective.  This isolates the precise set-theoretic collapse
  mechanism used to explain upper semi-compatibility;
\item the Remark 3.12.2(v) absorption distinction: zero-column unit
  indeterminacy is absorbed by a hull upper ray, while one-column logarithmic
  conjugacy is absorbed as log-volume equality; Lean now packages these two
  outcomes together with the corresponding Fig. 3.9 log-shell treatments, and
  instantiates the zero-column hull at the finite \(\Theta\)-boundary supplied
  by the Step (x) to Step (xi) corridor; Lean also records the exact
  inequality-versus-equality profile: two zero-column upper bounds and one
  one-column equality, hence inequalities in both directions on the one-column
  side; on the zero-column side, Lean proves that equality with the hull
  boundary, after identification with the finite \(\Theta\)-boundary from
  Step (xi), requires the corresponding reverse inequality as additional data;
\item the Remark 3.12.3(i) metric analogy:
  \(\operatorname{vol}(S)>0\) and \(\chi_S=-\operatorname{vol}(S)\) imply
  \(\chi_S<0\), hence \(\chi_S\le0\), together with the stored
  upper-semi-\((\mathrm{Ind3})\) analogue;
\item the Remark 3.12.4(iii) factor-\(p\) analogy: theta-label normalization
  divides a gap log-volume by \(\ell\), multiplication by \(\ell\) recovers it,
  and the positive factor \(\ell\) preserves nonpositive bounds exactly;
\item the Remark 3.12.4(v) p-adic analogue:
  an inclusion, not isomorphism, gives \((1-p)(2g-2)\le0\); under the
  recorded hypotheses \(p\ge2\) and \(g\ge2\), Lean also proves the defect is
  strictly negative;
\item the IUT IV, Theorem 1.10 algebraic handoff: an explicit expression for
  \(C_\Theta\), together with \(C_\Theta\ge-1\) and the remaining elementary
  estimates, implies the displayed \(\frac16\log(q)\) bound; Lean now exposes
  the exact Corollary 3.12 handoff identity
  \(C_\Theta+1=A-M\), where \(A\) is the arithmetic upper term and \(M\) is the
  main logarithmic term, so \(C_\Theta\ge-1\) gives \(0\le A-M\);
\item the final elementary coefficient estimates in IUT IV, Theorem 1.10:
  \((1-12/\ell^2)^{-1}\le2\) and
  \((1-12/\ell^2)^{-1}\frac14(1+36d_{\mathrm{mod}}/\ell)
   \le 1+80d_{\mathrm{mod}}/\ell\);
\item the IUT IV, Theorem 1.10 base-change monotonicity:
  \(\log(d_{F_{\mathrm{tpd}}})+\log(f_{F_{\mathrm{tpd}}})
   \le \log(d_F)+\log(f_F)\) propagates through the final right-hand side;
\item the IUT IV, Theorem 1.10 Step (i) identities:
  \(n^{-1}\sum_{m=1}^n m=(n+1)/2\) and
  \(n^{-1}\sum_{m=1}^n m^2=(2n+1)(n+1)/6\);
\item the IUT IV, Theorem 1.10 Step (ii) small-prime error:
  \(\log(2^{11}3^35^2)\le21\), represented by the corresponding linear
  logarithmic bound;
\item the IUT IV, Theorem 1.10 Step (i) \(\mathrm{GL}_2\) arithmetic constants
  for \(p=2,3,5\), including the product with the quadratic factor from (E5);
\item the IUT IV, Theorem 1.10 Step (iii) algebraic estimate combining the
  Step (ii) \(K\)-bound and the displayed \(\log(s_Q)\)-bound into
  \((1+4/\ell)\log(d_K,f_K)+(16/\ell)\log(s_Q)
   \le (1+36d_{\mathrm{mod}}/\ell)S+2\log(\ell)+70\);
\item the IUT IV, Theorem 1.10 Step (viii) Proposition 1.6 case split:
  the two cases \(e^*_{\mathrm{mod}}\ell\ge\eta_{\mathrm{prm}}\) and
  \(e^*_{\mathrm{mod}}\ell<\eta_{\mathrm{prm}}\) imply the uniform
  \(4/3(e^*_{\mathrm{mod}}\ell+\eta_{\mathrm{prm}})\) bound;
\item the IUT IV, Theorem 1.10 Step (viii) final absorption of
  \(2\log(\ell)+74\) and the prime-product term into
  \(10(e^*_{\mathrm{mod}}\ell+\eta_{\mathrm{prm}})\);
\item the IUT IV, Theorem 1.10 final display:
  the Corollary 3.12 \(C_\Theta\) handoff, followed by the
  tripodal-to-\(F\) monotonicity, gives the two displayed inequalities:
  first the \(F_{\mathrm{tpd}}\)-level \(\frac16\log(q)\) bound, then the
  final \(F\)-level bound;
\item the IUT IV, Theorem A bounded-discrepancy conclusion:
  boundedness below of
  \((1+\epsilon)(\log\operatorname{diff}_X+\log\operatorname{cond}_D)
   -\operatorname{ht}_{\omega_X(D)}\) is packaged together with the
  corresponding pointwise height inequality on points of bounded degree;
\item the IUT IV, Corollary 2.2(i) bounded-discrepancy chain:
  \(\frac16\log(q_2)\approx\frac16\log(q_\forall)
   \approx\frac16\operatorname{ht}_\infty
   \approx\operatorname{ht}_{\omega_X(D)}\) is packaged as three adjacent
  links together with the composed direct bounded-discrepancy relation;
\item bounded-discrepancy transfer lemmas: pointwise upper and lower bounds move
  across an \(\approx\) relation with the expected additive constant shift; Lean
  also packages simultaneous two-sided bound transfer;
\item the IUT IV, Corollary 2.2(ii) condition (C1) prime-scale window:
  \((\log(q_\forall))^{1/2}\le\ell\le
   10\delta(\log(q_\forall))^{1/2}\log(2\delta\log(q_\forall))\).  Lean keeps
  the square-root and logarithmic factor explicit, proves
  \(\log(q_\forall)\ge0\), and derives
  \(1\le10\delta\log(2\delta\log(q_\forall))\) from the nonempty window; Lean
  now packages the full C1 window endpoint with these immediate consequences;
\item the IUT IV, Corollary 2.2(ii) use of (C1) in Theorem 1.10:
  \(\sqrt h\le\ell\) and \(80d_{\mathrm{mod}}\le\delta\) imply
  \(80d_{\mathrm{mod}}/\ell\le\delta h^{-1/2}\); Lean packages this with the
  positive-denominator and nonnegative-numerator side conditions and the
  resulting \(1+\cdots\) coefficient inequality;
\item the IUT IV, Corollary 2.2(ii) error-term use of (C1):
  \(d^*_{\mathrm{mod}}\le\delta\) and
  \(\ell\le10\delta h^{1/2}\log(2\delta h)\) imply
  \(20(d^*_{\mathrm{mod}}\ell+\eta_{\mathrm{prm}})
   \le200\delta^2h^{1/2}\log(2\delta h)+20\eta_{\mathrm{prm}}\).  Lean
  packages this with the nonnegativity and upper-window premises needed for
  the multiplication comparison;
\item the IUT IV, Theorem 1.10 to Corollary 2.2(ii) first-bound composition:
  these two C1 replacements convert the Theorem 1.10 right-hand side into
  \((1+\delta h^{-1/2})S+
    200\delta^2h^{1/2}\log(2\delta h)+20\eta_{\mathrm{prm}}\);
\item the IUT IV, Corollary 2.2(ii) \(q_2/q\) comparison:
  the displayed estimate
  \(\frac16\log(q_2)-\frac16\log(q)
   \le\frac13h^{1/2}\log(2\delta h)\) follows through the intermediate
  term \(h^{1/2}\log(\ell)\), with \(1\le6\log(\ell)\) and
  \(3\log(\ell)\le\log(2\delta h)\) exposed in one endpoint;
\item the IUT IV, Corollary 2.2(ii) pre-\(\epsilon_E\) bound for
  \(h=\log(q_\forall)\):
  the Theorem 1.10 bound, the \(q_2/q\) estimate, and the bounded discrepancy
  \(B_K\) imply
  \[
    \frac16h\le(1+\delta h^{-1/2})S
    +(15\delta)^2h^{1/2}\log(2\delta h)+\frac12 C_K,
  \]
  with \(C_K=40\eta_{\mathrm{prm}}+2B_K\).  Lean exposes the
  \(q\to q_2\to q_\forall\) inequalities, the absorption into
  \((15\delta)^2\), and the identity \(\frac12C_K=20\eta_{\mathrm{prm}}+B_K\);
\item the IUT IV, Corollary 2.2(ii) condition (C2) inequality chain:
  \(\frac16\log(q)\le\frac16\log(q_2)\le\frac16\log(q_\forall)\le
   (1+\epsilon_E)(\log\operatorname{diff}_X+\log\operatorname{cond}_D)+C_K\);
\item the IUT IV, Corollary 2.2(ii) definition of \(\epsilon_E\):
  \(\epsilon_E=(60\delta)^2h^{-1/2}\log(2\delta h)\), represented with
  \(h^{1/2}\) and \(\log(2\delta h)\) as explicit real parameters; Lean proves
  positivity and \(\epsilon_E\ge5\delta h^{-1/2}\) from
  \(\delta\ge2\), \(h^{1/2}>0\), and \(\log(2\delta h)\ge1\), packaged as one
  definition endpoint;
\item the IUT IV, Corollary 2.2(ii) \(\epsilon_E\)-error conversion:
  substituting the definition of \(\epsilon_E\), Lean proves
  \((15\delta)^2h^{1/2}\log(2\delta h)
   \le \frac16h\cdot\frac25\epsilon_E\), packaged with
  \(h=(h^{1/2})^2\), \(h^{1/2}>0\), and the nonnegativity premise for the
  logarithmic factor;
\item the IUT IV, Corollary 2.2(ii) bridge to the denominator step:
  the pre-\(\epsilon_E\) \(h\)-bound, the error conversion, and
  \(\delta h^{-1/2}\le\epsilon_E/5\) yield the preliminary inequality used to
  move the \(\frac25\epsilon_E\frac16h\) term to the left; Lean packages this
  with log-degree nonnegativity and
  \(1+\delta h^{-1/2}\le1+\epsilon_E/5\);
\item the IUT IV, Corollary 2.2(ii) denominator step:
  Lean moves the term
  \((2\epsilon_E/5)\frac16h\) to the left and obtains the denominator
  \((1-2\epsilon_E/5)^{-1}\), packaged with positivity of
  \(1-2\epsilon_E/5\) and the preliminary inequality it acts on;
\item the IUT IV, Corollary 2.2(ii) final \(\epsilon_E\)-absorption:
  from \(0<\epsilon_E\le1\), Lean proves
  \[
    \frac{1+\epsilon_E/5}{1-2\epsilon_E/5}\le1+\epsilon_E,
    \qquad
    (1-2\epsilon_E/5)^{-1}C_K/2\le C_K,
  \]
  using nonnegativity of \(S\) and \(C_K\), and hence the displayed C2 bound
  from the preceding denominator expression;
\item the IUT IV, Corollary 2.2(ii) final \(h\)-bound:
  Lean composes the pre-\(\epsilon_E\) bound, error conversion, move-left step,
  and absorption to prove
  \(\frac16h\le(1+\epsilon_E)S+C_K\), with the square-root identity,
  \(\epsilon_E\) definition, coefficient comparison, \(0<\epsilon_E\le1\),
  and nonnegativity hypotheses exposed in one endpoint;
\item the IUT IV, Corollary 2.2(ii) final C2 bridge:
  using \(S\le\log\operatorname{diff}_X(x_E)+\log\operatorname{cond}_D(x_E)\),
  Lean derives the displayed C2 inequality chain, including the
  \(q\to q_2\to q_\forall\) inequalities and monotonicity under
  \(1+\epsilon_E\ge0\);
\item the IUT IV, Corollary 2.2(ii) comparison
  \(\log\operatorname{diff}_X(x_E)=\log(d_{F_{\mathrm{tpd}}})\) and
  \(\log(f_{F_{\mathrm{tpd}}})\le\log\operatorname{cond}_D(x_E)\le
   \log(f_{F_{\mathrm{tpd}}})+\log(2\ell)\); Lean packages these as the two
  corresponding sum inequalities for the tripodal and curve log sums;
\item the IUT IV, Corollary 2.2 to Theorem A passage:
  bounded-discrepancy equivalence
  \(\frac16\log(q_2)\approx\operatorname{ht}_{\omega_X(D)}\), together with the
  C2 bound for \(\frac16\log(q_2)\), gives a lower bound for
  \((1+\epsilon_E)(\log\operatorname{diff}_X+\log\operatorname{cond}_D)
   -\operatorname{ht}_{\omega_X(D)}\); Lean packages the lower
  bounded-discrepancy constant, the C2 input, the induced height upper bound,
  and the resulting discrepancy lower bound in one endpoint;
\item the IUT IV, Corollary 2.2/2.3 finite-exception passage:
  a lower bound on the complement of a finite exceptional set and a lower bound
  on the exceptional set glue to a global lower bound by taking the minimum of
  the two constants; Lean packages both minimum comparisons together with the
  resulting pointwise global lower bound;
\item the IUT IV, Corollary 2.3 Diophantine inequality:
  Lean records the GenEll transfer as an explicit hypothesis and derives the
  global bounded-discrepancy lower bound for
  \((1+\epsilon)(\log\operatorname{diff}_X+\log\operatorname{cond}_D)
   -\operatorname{ht}_{\omega_X(D)}\), packaged with \(d>0\),
  \(\epsilon>0\), the hyperbolic-curve premise, and the compact-subset
  Theorem A premise;
\item comparison levels:
  pointwise representative, aggregate representative, balanced sign-compatible,
  and hull/log-volume; Lean proves that the final weighted-theta route is
  hull/log-volume level and rejects the balanced sign-compatible level,
  including the balanced negation transport;
\item endpoint audits for the \(3.11.5 \Rightarrow 3.12\) comparison chart;
\item nonarchimedean and archimedean \((\mathrm{Ind3})\) local orientations;
\item ordered real-line alignment with cancellation only after scale equality;
\item a canonical unit-scale ordered alignment from a nonarchimedean
  \((\mathrm{Ind3})\) entry;
\item the local-object Hodge-descent square-weight boundary:
  transported zero-column and cusp-class local objects have the same finite
  log-volume as the packet object, while square/full-label preservation
  obligations remain explicit;
\item the square/full-label Hodge/SHE transport audit:
  a pointwise upper bound on the transported target full-label log-volumes
  transfers to the source square-weighted average once full-label log-volume,
  square weight, and total-weight preservation are supplied; the same transfer
  is now available inside the structured-SHE wrapper, and the resulting upper
  bound is exposed for the square-weighted average in the \(F_\ell\) cusp-label
  route after matching the source profile and source log-volume; consequently,
  structured-SHE target full-label bounds by the \(\Theta\)-source average feed
  the cusp-class container route all the way to
  \(q_{\mathrm{signed}}\le\Theta_{\mathrm{signed}}\); the same final comparison
  is now available directly from the primitive factored square/full-label SHE
  obligations, without requiring a prebuilt transport-audit argument;
\item a Gaussian-to-factored-SHE bridge:
  source and target degree-level Gaussian evaluations with equal environment
  degree supply the full-label value-preservation branch of the factored
  square/full-label SHE obligations; with separate coordinate-square and
  full-label-map preservation, Lean exposes the resulting pointwise comparison
  level, transported-average equality, forced identity of the representative
  square coordinate map, and retained separation of the two Hodge histories;
  the experiment route then reaches
  \(q_{\mathrm{signed}}\le\Theta_{\mathrm{signed}}\) directly from Gaussian
  target bounds; since the absolute-label exponent is nonnegative, Lean also
  derives those target bounds in the branch where the target environment degree
  is nonpositive and the \(\Theta\)-source average is nonnegative; conversely,
  the all-label Gaussian target-bound formulation forces the
  \(\Theta\)-source average to be nonnegative, since the core label has
  Gaussian degree \(0\); after removing the core label, Lean proves that every
  nonzero Gaussian target value is bounded by the \(\Theta\)-source average if
  the target environment degree is nonpositive and is itself bounded by that
  average; Lean also forms the corresponding finite average over the nonzero
  \(ZMod\) carrier, proves the same upper bound for this nonzero average, and
  proves the exact rescaling between this auxiliary nonzero average and the
  uniform \(F_\ell\)-average with zero Gaussian term;
  because the square-weighted route assigns weight \(0\) to the core label, Lean
  now feeds nonzero-only target bounds through the square/full-label transport
  audit all the way to \(q_{\mathrm{signed}}\le\Theta_{\mathrm{signed}}\); this
  nonzero Gaussian environment-comparison route is now a source-layer theorem,
  not only an experiment wrapper;
\item a bridge from factored square/full-label SHE obligations to the
  Hodge-descent \((\mathrm{Ind3})\) route.
\end{itemize}

\section{Current Lean Results}

The current first-pass experiment report evaluates to:
\[
\begin{array}{ll}
\text{route theorem flags} & = \texttt{true},\\
\text{collapse-obstruction flags} & = \texttt{true},\\
\text{Gaussian/procession finite-model flags} & = \texttt{true},\\
\texttt{balancedLevelRejectedAtFinalRouteTheoremAvailable} & = \texttt{true},\\
\texttt{disputeSettledByCurrentStage} & = \texttt{false}.
\end{array}
\]
Here the route flags are ordered real-line, nonarchimedean entry, and factored
SHE availability.  The obstruction flags are raw-cancellation failure,
label-independent \(j^2\) failure, and sign-quotient representative descent
failure.  The Gaussian/procession flags include finite tensor-packet, monoid
action, normalized splitting-generator average, the explicit square-sum formula
for the generator total and normalized delta, the resulting conditional
monotonicity of the acted tensor packet, canonical positive square-weight
profile, the constant-average check,
base-valuation product,
coric-transfer, determinant-normalization,
upper-ray, \(q\)-pilot two-computation, the Step (x) averaged-log-volume
indeterminacy corridor, the refined Step (x) bounds for the before-,
after-\((\mathrm{Ind1})\)-, and after-\((\mathrm{Ind2})\)-averages,
tensor-power warning, and global Frobenioid-calibration
shadows, plus the finiteness endpoint for
\(-|\log(\Theta)|\in\mathbb{R}\cup\{+\infty\}\), including extraction of the
finite real \(\Theta\)-value, the direct
Corollary 3.12 statement endpoint with \(\Theta\)-subject/\(q\)-not-subject
pilot bookkeeping, the Step (x)-sourced upper-ray, q-pilot two-computation,
statement endpoint, propagation of the same pilot-status boundary to that
endpoint, proof that the endpoint \(\Theta\)- and \(q\)-statuses are distinct,
identification of its \(q\)-pilot real with the Step (x)
pre-, after-\((\mathrm{Ind1})\)-, and after-\((\mathrm{Ind2})\)-averages,
positivity of \(|\log(q)|\) at that endpoint, and direct
\(-|\log(q)|\in\mathbb{R}_{\le-|\log(\Theta)|}\),
\(-|\log(q)|\le-|\log(\Theta)|\), finite-\(\Theta\) exclusion of \(+\infty\),
identification of the finite extended \(\Theta\)-value with the Step (x)
\((\mathrm{Ind3})\) upper bound, the nonnegative-\(\Theta\) opening reduction
including the stronger \(C_\Theta\ge0\) branch, and \(C_\Theta\nless-1\)
conclusions,
the
Step (xi-g)-to-(xi-f) endpoint bridge giving \(C_\Theta\ge-1\) from the
two-computation q-pilot record, the Step (xi-e)--(xi-f) fixed
\(-|\log(q)|\) membership check and its direct \(C_\Theta\ge-1\) consequence,
the composition into the statement-level endpoint, and positive-rational unit
rigidity.
The generalized-link flag records the rational parameter inequality
\(C_\Theta\ge-\lambda\), now with the corresponding \(q^\lambda\)-pilot
real-valued derivation.
The distinct-label flag records the logical intertwining display in
Remark 3.12.2, keeps the weakened abstract prime-strip condition fixed, and
rejects collapse of the q-native and \(\Theta\)-native labels; the final
Step (xi-f) inequality package is now checked under this same label separation.
The toy-model flag records the resulting \((\mathrm{ftoy})\) upper-ray endpoint
and arithmetic bound \(h\le\epsilon\).
The column-table flag records the Fig. 3.9 theta/q similarity-difference split:
upper similarities are common to both columns, while multiradiality and
log-shell treatment distinguish the 0-column from the 1-column; the upper
similarities are now exposed as a simultaneous zero/one-column theorem.
The absorption flag records the zero-column upper-bound versus one-column
equality distinction, tied to the two Fig. 3.9 log-shell treatments; the
one-column equality is now exposed together with its two induced inequalities,
while zero-column hull equality is equivalent to adding the missing reverse
bound; the same guard is now stated after replacing the hull by the finite
\(\Theta\)-boundary.
The metric flag records the Gauss--Bonnet sign calculation used as analogy for
the upper semi-compatibility inequality, including the nonpositive inequality
form following from strict negativity.
The log-link-juggling flag records the Remark 3.12.2(iii) source of
\((\mathrm{Ind3})\): unchanged procession-normalized averages through
\((\mathrm{Ind1})\) and \((\mathrm{Ind2})\), plus a one-sided upper-semi bound
and matching place families.
The factor-\(p\) flag records the \(F_\ell\)-label averaging denominator and
the order-preservation of this positive normalization.
The Frobenius-derivative flag records the degree inequality from inclusion and
the sharper strict negativity of the degree defect under \(p\ge2\), \(g\ge2\).
The IUT IV flag records the downstream Theorem 1.10 use of the Corollary 3.12
lower bound on \(C_\Theta\), including the identity \(C_\Theta+1=A-M\), the
resulting nonnegative gap \(0\le A-M\), the final coefficient estimates, and the
tripodal-to-\(F\) log-degree monotonicity.  The Step (i) flag records the two
elementary sum identities used in the leading-term computation.  The Step (ii)
small-prime flag records the additive error bound \(21\).  The \(\mathrm{GL}_2\)
flag records the finite cardinality constants feeding that bound.  The
Step (iii) flag records the \(\log(s_Q)\) coefficient absorption.
The Step (viii) flag records the prime-product case split used before the final
upper bound for \(C_\Theta\), plus the final coarse error absorption.  The final
display flag records the composed Theorem 1.10 inequality with \(F\)-level
log-degree terms.  The Theorem A flag records the bounded-discrepancy
height/log-different/log-conductor inequality.  The Corollary 2.2 flag records
composition of the displayed bounded-discrepancy chain.  The transfer flag
records movement of pointwise bounds across bounded-discrepancy equivalence.
The Corollary 2.2(C1) flag records the displayed prime-scale window and its
elementary consequences.  The Corollary 2.2(C2) flag records the displayed
three-step logarithmic inequality chain.  The C1 coefficient flag records the
replacement of \(80d_{\mathrm{mod}}/\ell\) by \(\delta h^{-1/2}\) in the
Theorem 1.10-to-C2 proof.  The C1 error-term flag records the corresponding
replacement of \(20d^*_{\mathrm{mod}}\ell\) by
\(200\delta^2h^{1/2}\log(2\delta h)\).  The first-bound flag records the
composition of those replacements with Theorem 1.10.  The \(\epsilon_E\)-definition flag
records the source formula and lower estimate for \(\epsilon_E\).  The
\(\epsilon_E\)-move-left flag records introduction of the denominator
\((1-2\epsilon_E/5)^{-1}\).  The \(\epsilon_E\)-absorption flag records the
final denominator removal in C2.  The
\(\epsilon_E\)-error flag records the conversion of the \((15\delta)^2\)-term
to the \(\frac25\epsilon_E\frac16h\) term.
The final \(h\)-bound flag records the full composition through absorption.
The final C2 bridge flag records the last monotonic replacement of the tripodal
log-degree sum by the curve different/conductor sum.
The \(h\)-to-move-left flag records the bridge from the pre-\(\epsilon_E\) bound
to the denominator-stage preliminary inequality.
q2/q flag records the displayed comparison of \(\log(q_2)\) and \(\log(q)\).
The pre-\(\epsilon_E\) \(h\)-bound flag records the subsequent absorption of
the bounded \(q_\forall/q_2\) discrepancy into \(C_K\).
The log-different/log-conductor flag records the bridge from tripodal field
terms to curve functions.  The Corollary 2.2-to-Theorem A flag records the formal
bounded-discrepancy transfer from C2 to the Theorem A inequality.
The finite-exception flags record the lower-bound gluing used when finite
exceptional sets are enlarged in Corollary 2.2 and then passed toward
Corollary 2.3.
The Corollary 2.3 flag records the final GenEll-transfer interface for the
Diophantine inequality.
The theorem-level route is:
\[
\begin{array}{ccccc}
\text{nonarch. }(\mathrm{Ind3})\text{ entry}
& \to &
\text{canonical ordered } \R\text{-alignment}
& \to &
\text{source/target alignment}
\\[3pt]
& & \downarrow & & \downarrow
\\[-2pt]
\text{factored SHE transport}
& \to &
\text{Hodge-descent audit}
& \to &
q_{\mathrm{signed}}\le \Theta_{\mathrm{signed}}.
\end{array}
\]
The route remains conditional on the displayed hypotheses.
The local-object Hodge-descent layer now exposes equality of the transported
finite log-volumes before square/full-label preservation is supplied.
The square/full-label audit now also transports target pointwise upper bounds
back to the source square-weighted average, including through the structured-SHE
route wrapper, and the experiment layer exposes the resulting final
q-to-\(\Theta\) comparison theorem.  The same experiment-level theorem is now
exposed for primitive factored square/full-label SHE obligations and for
degree-level Gaussian evaluations after the explicit preservation branches are
supplied.  In the nonpositive-environment/nonnegative-\(\Theta\)-average branch,
the Gaussian pointwise target bound is derived from the sign of the
\(q^{j^2}\)-exponent rather than assumed.  Lean also proves that this
all-label Gaussian bound cannot hold when the \(\Theta\)-source average is
negative; under nonpositive environment it is equivalent to nonnegativity of
that average.  Lean also packages the resulting separation: for negative
\(\Theta\)-source average, the all-label bound fails, while the nonzero-label
bound still holds whenever the Gaussian environment degree is bounded by that
same average.  With the remaining endpoint hypotheses, this same negative
\(\Theta\)-branch still yields the final \(q_{\mathrm{signed}}\le
\Theta_{\mathrm{signed}}\) route through the nonzero-bound endpoint; the
target nonpositive-environment hypothesis is then available, but the
nonnegative-\(\Theta\) hypothesis of the nonpositive/nonnegative all-label
endpoint is impossible.  The
source-supplied nonzero/environment formulation gives the same conclusion and
rejects both source and target all-label bounds.  More
generally, Lean decomposes the all-label Gaussian bound
exactly into nonnegativity of the upper bound at the zero label and the
corresponding nonzero-label bound.  The corresponding identity-coordinate final
route can now be stated directly from canonical \(j=1\) equality, nonpositive
target environment, and nonnegative \(\Theta\)-source average; if the all-label
Gaussian target bound is supplied directly, the same endpoint no longer needs
the nonpositive/nonnegative branch hypotheses, and its proof also uses only the
nonzero-label part of the supplied all-label target bound.  The nonzero environment route
also has a source-side variant: canonical \(j=1\) equality transports
nonpositivity and the comparison with the \(\Theta\)-source average from the
source environment degree to the target environment degree.  Lean also proves
that the nonzero target-bound predicate is then equivalent to this
source-environment comparison, hence to the source nonzero-bound predicate
under source nonpositivity, and that the all-label predicate is equivalent
to this source comparison plus nonnegativity at the zero label.  The final
all-label q-to-\(\Theta\) endpoint now consumes this source-side package
directly through the source-environment comparison; the zero-label
nonnegativity component is retained in the API but not used by this final
proof.  It can alternatively derive that package from a source all-label
Gaussian bound under source nonpositivity; the proof now uses only the
nonzero-label part of this source all-label input.  The source all-label
predicate is itself equivalent, under source nonpositivity, to the conjunction of
\(0\le\Theta_{\mathrm{avg}}\) and the source-environment upper bound.  Hence,
under canonical \(j=1\) preservation and the two nonpositive branches, the
source and target all-label Gaussian predicates are equivalent.  It can also
consume the decomposed zero/nonzero input directly; this endpoint now ignores
the zero-label nonnegativity component and proceeds through the nonzero
factored route.  In the negative \(\Theta\)-branch, Lean proves that this
zero-label nonnegativity component is impossible, while the source
environment/nonzero route still gives the final comparison.  For nonzero labels, Lean
proves the complementary bound: the absolute exponent is at least \(1\), so a
nonpositive environment degree below the \(\Theta\)-source average bounds every
nonzero Gaussian component and hence their finite nonzero-label average.  The
final nonzero route may now consume the source nonzero Gaussian predicate
directly after canonical \(j=1\) preservation and source nonpositivity; Lean
reduces this input to the source-environment upper bound before invoking the
source-environment endpoint.  The
square-weighted transport route has also been strengthened from an all-label
pointwise bound to a weighted-summand bound, so the zero label requires no
target upper bound.  The experiment theorem now delegates to the source-layer
q-to-\(\Theta\) theorem for this nonzero Gaussian environment-comparison route.
Separately, Lean proves that the
uniform \(F_\ell\)-average of the Gaussian shadow is
\((\ell-1)/\ell\) times the auxiliary nonzero-carrier average, fixing the
denominator at \(\ell\) when the zero label is present.

\section{Audit Findings}

No formal statement currently proves IUT III, Corollary 3.12 from Mochizuki's
published hypotheses.  The strongest positive result is conditional: if the
hull/log-volume upper-ray comparison supplies
\(-|\log(q)|\le-|\log(\Theta)|\), then Lean checks the formal Step (xi-f)
conclusion \(C_\Theta\ge-1\); if the upper-ray membership is strict, Lean checks
the sharper \(C_\Theta>-1\), while the boundary value \(C_\Theta=-1\) collapses
the q-pilot and \(\Theta\)-boundary log-volumes to equality and forces the
\(\Theta\)-boundary to be negative.  If the
nonarchimedean \((\mathrm{Ind3})\) entry,
ordered real-line alignment, and Hodge/SHE square-weight transport audit are
supplied, Lean checks the final signed \(q\)-to-\(\Theta\) inequality feeding
that algebra.  If the corresponding primitive factored square/full-label SHE
obligations are supplied instead, Lean first constructs the transport audit and
checks the same final inequality.  If the input is a pair of degree-level
Gaussian evaluations with equal environment degree, Lean constructs the
full-label log-volume compatibility branch before invoking the factored route.
The branch with nonnegative \(\Theta\)-source average is formally easier:
nonnegativity of \(j^2\) makes every nonpositive environment-degree Gaussian
target value bounded above by that average.  The all-label version cannot be
the negative branch: the core label contributes \(0\), so a negative upper bound
is impossible; Lean now packages this as an exact equivalence with
nonnegativity of the upper bound and as a separation theorem from the
nonzero-label bound under the environment comparison; the separated
nonzero-label branch also feeds the final \(q\)-to-\(\Theta\) theorem, both
from target-side and source-side environment data.  In the source-side version,
Lean rejects both all-label predicates while retaining the final comparison.
For the target-side nonpositive/nonnegative endpoint, Lean isolates the failed
premise as \(0\le\Theta_{\mathrm{avg}}\), not target nonpositivity.
On the nonzero-label part, the formal
obstruction disappears:
exponent \(\ge1\) gives the expected bound from the environment comparison.  The
square-weighted final route now accepts this nonzero-only bound, since the core
summand has weight \(0\).  This gives a conditional source-level q-to-\(\Theta\)
route from degree-level Gaussian data, coordinate-square preservation,
full-label-map preservation, equal source/target environment degree, and
environment degree bounded above by the \(\Theta\)-source average.  In the
identity-coordinate case, Lean now supplies the coordinate-square and
full-label-map preservation branches internally, so this route only asks for
equal Gaussian environment degree and the nonpositive environment comparison;
equivalently, the equal-environment input may be replaced by equality of the
canonical \(j=1\) Gaussian components, and the target-side environment
comparison may be supplied on the source side after the same equality; the
nonzero target-bound predicate is equivalent to this source-side comparison,
and the all-label target predicate is equivalent to this comparison together
with \(0\le\Theta_{\mathrm{avg}}\).  The corresponding final theorem now
uses these source-side inputs directly rather than asking for the all-label
target predicate separately; a further variant derives them from a source
all-label Gaussian bound, and Lean classifies that source all-label predicate
as exactly \(0\le\Theta_{\mathrm{avg}}\) plus the source-environment bound; the
target and source all-label predicates are consequently equivalent under the
same canonical \(j=1\) and nonpositivity hypotheses.
Lean also proves that, in the nonpositive
environment branch, the nonzero coordinate target-bound formulation is
equivalent to this environment bound against the
\(\Theta\)-source average and, with source nonpositivity, to the source
nonzero predicate;
the necessity direction does not use the nonpositive hypothesis, since the
canonical \(j=1\) component is the environment degree itself.
The final identity-coordinate route may now be stated directly from canonical
\(j=1\) equality and the nonzero Gaussian target-bound predicate; after source
nonpositivity, the corresponding source nonzero-bound predicate supplies the
same endpoint through the source-environment classification.
In the all-label branch, the analogous direct endpoint may either use canonical
\(j=1\) equality plus the all-label Gaussian target-bound predicate, or derive
that predicate from nonpositive target environment and nonnegative
\(\Theta\)-source average.  Lean also records the inclusion from the all-label
predicate to the nonzero-label predicate used by the weighted-summand route,
and the converse after adding the zero-label nonnegativity condition; the
corresponding direct endpoint is exposed with this decomposed input.  The
target all-label endpoint is now implemented by this same nonzero/square-weighted
route rather than by a full all-label factored bound.
The
zero/nonzero average split is now exact: the nonzero carrier has cardinality
\(\ell-1=2j_{\max}\), the signed carrier is a twofold cover of the nonzero
absolute half-range, the zero Gaussian degree is \(0\), and the all-\(F_\ell\)
average is \((\ell-1)/\ell\) times the nonzero absolute average; consequently
the \(F_\ell\)-coordinate and full \(|F_\ell|\)-absolute averages differ
strictly for nonzero environment degree.  Lean also now proves the local
distinct-label fact emphasized in IUT II: the \(j=1\) and \(j=2\) theta/Gaussian
components coincide if and only if the environment degree vanishes.  The
remaining task is to replace
the current degree-level Gaussian shadow by the actual Hodge--Arakelov
Gaussian/Frobenioid construction behind this environment comparison.

The strongest negative result is also limited: in the current \(\F_\ell\)-label
model, Lean rejects a single label-independent scale accounting for all
representative \(j^2\) factors, and it rejects balanced sign-compatible evidence
as the final hull/log-volume comparison level.  Lean also rejects direct descent
of arbitrary raw representative \(j^2\) theta-degree data to \(|F_\ell|\), while
the underlying \(F_\ell\)-label model, additive torsor laws, sign-unit quotient,
and canonical cusp-class bridge are now packaged as source-side endpoints,
accepting the half-range exponent convention stated in IUT II.
Lean also checks that an \(N\)-th tensor-power boundary is a strict strengthening
for \(N\ge2\) and negative pilot log-volume.  It also checks the Step (xii)
numerical warning: nonzero local Frobenioid shifts change log-volume, whereas
global calibration restores the unshifted value.  The positive rational unit
rigidity fact in Remark 3.12.1(iii) is formalized directly over \(\mathbb{Q}\).
For generalized links, Lean records that \(\lambda<1\) makes \(-\lambda>-1\).
At \(\lambda=1\), the generalized endpoint reduces to the ordinary q-pilot
boundary.
Lean also checks that the Remark 3.12.2 simultaneous-intertwining step depends
on distinct labels rather than an unlabeled identification, while the weakened
prime-strip condition is retained.
In the accompanying toy model, Lean proves the final upper bound exactly from
the upper-ray membership hypothesis.
Lean also records that the q-pilot column lacks the theta-side multiradiality
used in Theorem 3.11(iii)(c).
For Remark 3.12.2(iv), Lean records that the essential \(\bullet_0\) and
\(\circ\) roles, log-volume compatibility, and log-Kummer non-interference hold
on both the 0-column and the 1-column.
For Remark 3.12.2(iii), Lean packages the log-link juggling as the source of
the upper-semi \((\mathrm{Ind3})\) endpoint: the averaged value is fixed through
\((\mathrm{Ind1})\) and \((\mathrm{Ind2})\), then bounded from above.
For Remark 3.12.2(v), Lean separates hull absorption of unit shifts from
log-volume equality after logarithmic-conjugate compatibility, and records that
q-pilot non-interference does not supply q-pilot multiradiality.  It also checks
that zero-column hull upper bounds do not become hull equalities without
reverse inequalities, even after the hull is identified with the finite
\(\Theta\)-boundary in Step (xi).
For Remark 3.12.3(i), Lean proves the stated negative Euler-characteristic sign
from positivity of metric volume.
For Remark 3.12.4(iii), Lean checks the \(\ell\)-denominator normalization
identity for the theta-label averaging analogy.
For Remark 3.12.4(v), Lean proves the displayed nonpositive degree defect from
\(p\ge2\) and \(g\ge2\).
For IUT IV, Theorem 1.10, Lean proves the algebraic implication from the
computed \(C_\Theta\) expression and \(C_\Theta\ge-1\) to the final
\(\frac16\log(q)\) estimate, conditional on the separate elementary estimates.
Lean now also proves the two displayed final elementary coefficient estimates
from \(\ell\ge5\).
It also proves the final monotonic replacement of the tripodal log-degree sum by
the larger \(F\)-level log-degree sum, conditional on the two component
monotonicity inequalities cited from Proposition 1.3(i).
Lean also proves the Step (i) sum and square-sum identities over \(\mathbb{R}\).
For Step (ii), Lean proves the small-prime additive error bound from the
displayed logarithmic inequalities.
Lean also checks the Step (i) \(\mathrm{GL}_2\) constants for \(p=2,3,5\) and
their product with the quadratic extension factor.
For Step (iii), Lean proves the displayed coefficient absorption from the
source hypotheses \(d_{\mathrm{mod}}\ge1\), \(\ell\ge5\), \(2\log(\ell)\le\ell\),
and nonnegativity of the tripodal log-degree sum.
For Step (viii), Lean checks that the two Proposition 1.6 cases imply the
uniform prime-product bound used in the final \(C_\Theta\) estimate.
Lean also checks the subsequent absorption into the coefficient \(10\) of
\(e^*_{\mathrm{mod}}\ell+\eta_{\mathrm{prm}}\).
Combining the Corollary 3.12 handoff with the tripodal-to-\(F\) monotonicity,
Lean proves the two-line final display of Theorem 1.10 at this real-valued
shadow level: the intermediate \(F_{\mathrm{tpd}}\)-level bound and the final
\(F\)-level bound.
For Theorem A, Lean records the exact bounded-discrepancy conclusion and
packages it with the equivalent pointwise height bound from the lower-bound
hypothesis.
For Corollary 2.2(i), Lean defines equality up to bounded discrepancy, proves
reflexivity, symmetry, transitivity, nonnegative scaling, and composes the
displayed chain from \(\log(q_2)\) to \(\operatorname{ht}_{\omega_X(D)}\);
the chain endpoint records both the adjacent links and the composite relation.
Lean also proves that upper and lower pointwise bounds transfer across such a
bounded-discrepancy equivalence with the appropriate constant shift, including a
combined two-sided transfer theorem.
For Corollary 2.2(ii)(C1), Lean records the displayed window for \(\ell\) and
derives nonnegativity of \(\log(q_\forall)\), positivity of
\((\log(q_\forall))^{1/2}\), and
\(1\le10\delta\log(2\delta\log(q_\forall))\), packaged as a single C1 endpoint.
Lean also verifies the C1 coefficient replacement
\(80d_{\mathrm{mod}}/\ell\le\delta h^{-1/2}\) used immediately after applying
Theorem 1.10, together with the positivity/nonnegativity facts needed for the
division comparison.
Lean also verifies the C1 upper-window replacement for the Theorem 1.10 error
term \(20(d^*_{\mathrm{mod}}\ell+\eta_{\mathrm{prm}})\), including the
nonnegativity and upper-window facts used in the product estimate.
Combining these, Lean proves the first displayed post-Theorem-1.10 bound in the
proof of Corollary 2.2(ii).
Lean also proves the displayed comparison
\(\frac16\log(q_2)-\frac16\log(q)
\le\frac13h^{1/2}\log(2\delta h)\) from the intermediate
\(h^{1/2}\log(\ell)\) bounds and their two logarithmic side inequalities.
Lean then combines this with the bounded \(q_\forall/q_2\) discrepancy to prove
the displayed pre-\(\epsilon_E\) bound for \(\frac16h\), with the
\(q\to q_2\to q_\forall\) transfers, \((15\delta)^2\) absorption, and
\(\frac12C_K\) rewrite exposed as named subclaims.
For Corollary 2.2(ii), Lean proves the final \(\frac16\log(q)\) bound from the
displayed \(\log(q)\), \(\log(q_2)\), and \(\log(q_\forall)\) chain.
Lean records the proof's definition of \(\epsilon_E\) and verifies the stated
lower estimate \(\epsilon_E\ge5\delta h^{-1/2}\), including the positivity
facts used in the inverse-square-root comparison.
Lean also verifies the source conversion
\((15\delta)^2h^{1/2}\log(2\delta h)\le\frac16h\cdot\frac25\epsilon_E\),
with \(h=(h^{1/2})^2\) and the required nonnegativity facts explicit.
Combining this with \(\delta h^{-1/2}\le\epsilon_E/5\), Lean proves the
preliminary inequality needed for the move-left step, including the
nonnegative log-degree premise used to scale the coefficient comparison.
Lean then verifies the algebraic move-left step that introduces the denominator
\((1-2\epsilon_E/5)^{-1}\), exposing denominator positivity and the input
preliminary inequality.
Lean also proves the final \(\epsilon_E\)-absorption step used just before C2:
the intermediate expression with denominator \(1-2\epsilon_E/5\) is bounded by
\((1+\epsilon_E)\) times the log-degree sum plus \(C_K\), with denominator
\(\ge1/2\) and nonnegativity of \(S\) and \(C_K\) explicit.
Lean now composes these steps into the final displayed \(h=\log(q_\forall)\)
bound, with all square-root, \(\epsilon_E\), coefficient, and nonnegativity
premises exposed in one endpoint.
Lean also derives the displayed C2 chain from this \(h\)-bound and the
tripodal-to-curve log-degree comparison, recording the \(q\to q_2\to
q_\forall\) inequalities and the monotonicity condition \(1+\epsilon_E\ge0\).
It also proves the two sum inequalities relating the tripodal different and
conductor terms to \(\log\operatorname{diff}_X+\log\operatorname{cond}_D\),
packaged with the defining equality and conductor inequalities.
For the Corollary 2.2 to Theorem A passage, Lean combines the
bounded-discrepancy lower estimate for
\(\frac16\log(q_2)-\operatorname{ht}_{\omega_X(D)}\) with the C2 upper bound
for \(\frac16\log(q_2)\), deriving both the corresponding upper bound for the
height and the lower bound for
\((1+\epsilon_E)(\log\operatorname{diff}_X+\log\operatorname{cond}_D)
-\operatorname{ht}_{\omega_X(D)}\).
For the Corollary 2.2/2.3 exceptional-set passage, Lean proves that a lower
bound off a finite exceptional set and a lower bound on the exceptional set imply
a global lower bound, and specializes this to the Theorem A discrepancy.  The
endpoint records both minimum comparisons and the resulting pointwise bound.
For Corollary 2.3, Lean proves the stated bounded-discrepancy conclusion from an
explicit GenEll transfer hypothesis; GenEll itself is not formalized here.  The
endpoint exposes \(d>0\), \(\epsilon>0\), hyperbolicity, and the compact-subset
Theorem A premise.
Finally, Lean combines the Corollary 2.2(i) bounded-discrepancy height
comparison with the C2 bound to construct the Theorem A lower-bound shadow.

\section{Missing Mathematics}

The present code does not yet formalize:
\begin{itemize}
\item initial \(\Theta\)-data as in IUT I;
\item actual Hodge theaters, Frobenioids, prime strips, log-shells, or
  \(\Theta\)-links;
\item the Hodge--Arakelov theta evaluation and Gaussian monoids of IUT II;
\item the construction of the actual theta values, Gaussian/splitting monoids,
  \(F_{\mathrm{mod}}^\times\)-actions, Kummer isomorphisms,
  mono-analyticization, and Frobenioids behind the present shadows;
\item the actual construction of log-shells and tensor packets attached to
  mono-analytic prime-strips in IUT III, Theorem 3.11;
\item the genuine holomorphic hull and determinant operations of Step (xi),
  beyond the present finite determinant/upper-ray/two-computation/tensor-power
  shadows;
\item the global realified Frobenioids of Step (xii), beyond the present integer
  shift and calibrated-\(C_\Theta\) shadow;
\item the arithmetic-divisor and analytic local log-volume construction behind
  the IUT IV Theorem 1.10 numerical estimates, beyond the present source-side
  interface records;
\item the construction of the normalized height, log-different, and log-conductor
  functions used in IUT IV, Theorem A;
\item the construction of the bounded-discrepancy equivalences asserted in
  IUT IV, Corollary 2.2(i);
\item a proof that the current Hodge/SHE transport obligations follow from
  Mochizuki's full hypotheses.
\end{itemize}

\section{Next Work}

The finite \(q^{j^2}\), splitting-generator, and Gaussian-to-factored-SHE
packaging endpoints are now in place.  The next significant formal step is to
replace current explicit Hodge/SHE transport assumptions by constructions
closer to the papers:

\[
\begin{array}{c}
\text{theta values }q^{j^2}\quad\checkmark\ \text{finite square-weight endpoint}\\
\leadsto\ \text{Gaussian and splitting monoids}\quad\checkmark\
  \text{generator action endpoint}\\
\leadsto\ \text{factored square/full-label preservation}\quad\checkmark\
  \text{Gaussian endpoint}\\
\leadsto\ \text{hull/log-volume route}.
\end{array}
\]
The next step is to construct the Hodge/SHE preservation hypotheses themselves.

\begingroup
\small
\begin{thebibliography}{9}
\bibitem{IUT1}
S. Mochizuki, \emph{Inter-universal Teichmuller Theory I:
Construction of Hodge Theaters}.

\bibitem{IUT2}
S. Mochizuki, \emph{Inter-universal Teichmuller Theory II:
Hodge-Arakelov-theoretic Evaluation}.

\bibitem{IUT3}
S. Mochizuki, \emph{Inter-universal Teichmuller Theory III:
Canonical Splittings of the Log-theta-lattice}.

\bibitem{IUT4}
S. Mochizuki, \emph{Inter-universal Teichmuller Theory IV:
Log-volume Computations and Set-theoretic Foundations}.

\bibitem{MochizukiLean}
S. Mochizuki, \emph{On the Formalization of IUT: Preliminary Progress Report},
2026.

\bibitem{ScholzeStix}
P. Scholze and J. Stix, \emph{Why abc is still a conjecture}.
\end{thebibliography}
\endgroup

\end{document}
